\chapter{Eksik Veri ve Veri Kalitesi}
\label{ch:missing-data-quality}
\index{eksik veri}
\index{veri kalitesi}
\index{veri temizleme}

\begin{hedefler}
Bu bölümün sonunda okuyucu:
\begin{itemize}
  \item gerçek sıfır, geçersiz kod, yapısal eksiklik ve yanıtlanmamış değeri ayırt edebilecek,
  \item eksik gözlem oranını değişken, kişi, grup ve örüntü düzeyinde özetleyebilecek,
  \item MCAR, MAR ve MNAR mekanizmalarının varsayımlarını ve sınırlarını açıklayabilecek,
  \item tam gözlem, mevcut çift, tekli atama, model temelli yöntem ve çoklu atamayı karşılaştırabilecek,
  \item eksik veri kararının tahmin, standart hata ve genellenebilirlik üzerindeki etkisini değerlendirebilecek,
  \item duyarlılık analizi ile ana sonucun eksik veri varsayımına bağımlılığını inceleyebilecek,
  \item Python, SPSS ve R ile veri kalite denetimi yapabilecek,
  \item eksik veri yöntemini şeffaf ve yeniden üretilebilir biçimde raporlayabilecektir.
\end{itemize}
\end{hedefler}

\begin{hazirlik}
Bir öğrenci veri setinde sınav puanı için boş hücre, 0, 99 ve ``girmedi''
değerleri görülmektedir. Bu dört gösterimin neden doğrudan aynı işleme tabi
tutulamayacağını açıklayın.
\end{hazirlik}

\section{PDR vakası: İzlem ölçümünü kimler tamamlamadı?}

Bir psikolojik danışma merkezi, altı haftalık sınav kaygısı programının
sonunda öğrencilerin kaygı puanlarını yeniden ölçmektedir. Programa başlayan
120 öğrencinin 18'i son ölçümü tamamlamamıştır. Yalnız ``102 geçerli, 18 eksik''
demek yeterli değildir. İzlemi tamamlamayan öğrenciler başlangıçta daha yüksek
kaygı puanına sahipse veya programdan yarar görmediği için ayrılmışsa, kalan
102 öğrencinin ortalaması program sonucunu gereğinden olumlu gösterebilir.

Eksik veri yalnız hücre sayısını azaltmaz. Hangi öğrencilerin, hangi
değişkenlerde ve hangi nedenle eksik kaldığı; hedeflenen bilimsel sonucu ve
uygun analiz yöntemini değiştirebilir \cite{little2019,rubin1976}.

\begin{researchbox}
Eksiklik oranı küçük olsa bile kritik bir alt grupta yoğunlaşabilir. Yüzde 5
eksiklik bütün gruplara dağılmışsa etkisi sınırlı olabilir; aynı yüzde 5 yalnız
yüksek risk grubunda oluşmuşsa hizmet gereksinimi tahmini ciddi biçimde
yanlılaşabilir.
\end{researchbox}

\section{Önce veri kalitesini tanımlamak}
\index{veri kalitesi!boyutları}

Veri kalitesi yalnız eksik hücre bulunmaması değildir. Analizden önce en az şu
boyutlar incelenmelidir:

\begin{description}
  \item[Doğruluk] Değer özgün kayıt ve ölçme süreciyle uyuşuyor mu?
  \item[Tamlık] Gerekli değişken ve zaman noktalarının ne kadarı gözlenmiş?
  \item[Tutarlılık] Birbiriyle ilişkili alanlar mantıksal olarak uyumlu mu?
  \item[Geçerlilik] Değer, veri sözlüğündeki tür, aralık ve kodlara uyuyor mu?
  \item[Benzersizlik] Aynı analiz birimi yanlışlıkla birden çok kez kaydedilmiş mi?
  \item[Güncellik] Ölçüm zamanı araştırma sorusunun dönemini temsil ediyor mu?
\end{description}

Örneğin yaşın 250 olması geçersiz değer, aynı öğrencinin iki kez kaydedilmesi
yinelenen kayıt, son-test tarihinin programdan önce olması zaman tutarsızlığı,
boş son-test hücresi ise eksik veridir. Her sorun farklı bir düzeltme ve kayıt
gerektirir.

\begin{definitionbox}
\textbf{Eksik veri}, araştırmada tanımlanan bir değişkenin belirli bir analiz
birimi için gözlenmemiş olmasıdır. Eksiklik; değerin sıfır, düşük, olumsuz veya
uygulanamaz olduğu anlamına gelmez.
\end{definitionbox}

\section{Eksik değer, kod hatası ve yapısal eksiklik}
\index{yapısal eksiklik}
\index{eksik değer kodu}

Boş hücrelerin tümü aynı anlamı taşımaz. Tablo~\ref{tab:missing-types}, yaygın
durumları ayırır.

\begin{table}[htbp]
\centering
\caption{Eksiklik ve veri kalitesi durumlarının ayrımı.}
\label{tab:missing-types}
\begin{tabular}{@{}p{0.21\textwidth}p{0.33\textwidth}p{0.30\textwidth}@{}}
\toprule
Durum & Örnek & Uygun ilk işlem \\
\midrule
Yanıtlanmamış değer & Öğrenci son-test formunu doldurmamış & Eksik olarak tanımla; neden ve örüntüyü incele \\
Yapısal eksiklik & Çocuğu olmayan kişiye ebeveynlik sorusu uygulanmamış & ``Uygulanamaz'' durumunu gerçek eksiklikten ayır \\
Geçersiz kod & 1--5 maddesinde 9 değeri & Özgün formu kontrol et; doğrulanamıyorsa eksik yap \\
Gerçek sıfır & Öğrencinin devamsızlık günü 0 & Geçerli değer olarak koru \\
Ölçüm sınırı & Süre cihazın alt sınırının altında & Sansürlü/ölçüm sınırlı veri olarak belgele \\
\bottomrule
\end{tabular}
\end{table}

``99'', ``999'' veya $-1$ gibi özel kodlar analize girmeden önce yazılımın
gerçek eksik değer gösterimine dönüştürülmelidir. Ancak bu dönüşüm özgün
değişken üzerine yazılmamalı; ham veri korunmalı ve temizleme betiğinde yeni
analiz değişkeni üretilmelidir.

\begin{mistakebox}
Eksik puanı 0 yapmak, ``yanıt yok'' bilgisini gerçek ve en düşük puan gibi
yorumlar. Bu işlem ortalamayı, varyansı, korelasyonu ve grup farklarını
sistematik biçimde değiştirebilir.
\end{mistakebox}

\section{Eksikliğin miktarı ve örüntüsü}
\index{eksik veri örüntüsü}

Tek bir genel eksiklik yüzdesi yerine dört düzey birlikte incelenir:

\begin{enumerate}
  \item her değişkendeki eksik sayı ve oran,
  \item her kişideki eksik değişken sayısı,
  \item gruplara ve önemli başlangıç özelliklerine göre eksiklik oranı,
  \item değişkenlerin birlikte eksik kaldığı örüntüler.
\end{enumerate}

Bir değişken için eksiklik göstergesi
\[
R_i=
\begin{cases}
1, & \text{değer gözlenmişse},\\
0, & \text{değer eksikse}
\end{cases}
\]
biçiminde tanımlanabilir. $R$ göstergesinin grup, başlangıç puanı veya başka
gözlenen değişkenlerle ilişkisi, eksikliğin rastgele dağılmadığını gösterebilir.
Bu inceleme eksiklik mekanizmasını kesin olarak kanıtlamaz; yalnız makul ve
makul olmayan varsayımları ayırmaya yardım eder.

\section{MCAR, MAR ve MNAR mekanizmaları}
\index{MCAR}
\index{MAR}
\index{MNAR}

Eksik veri mekanizmaları, eksik kalma olasılığının hangi bilgilere bağlı
olduğunu ifade eder \cite{rubin1976,little2019,vanbuuren2018}.

\begin{description}
  \item[MCAR: tamamen rassal eksiklik] Eksik kalma olasılığı gözlenen ve
  gözlenmeyen değerlerle ilişkili değildir. Örneğin bazı formların taşıma
  sırasında rastgele kaybolması bu varsayıma yaklaşabilir.
  \item[MAR: gözlenen bilgiye koşullu rassal eksiklik] Eksiklik, gözlenen
  değişkenlerle açıklanabilir; aynı gözlenen bilgilere sahip kişiler arasında
  eksiklik gözlenmeyen değere ayrıca bağlı değildir. Son-test eksikliğinin
  gözlenen başlangıç puanı ve program grubuyla açıklanması buna örnektir.
  \item[MNAR: rassal olmayan eksiklik] Gözlenen değişkenler hesaba katıldıktan
  sonra bile eksiklik gözlenmeyen değerin kendisiyle ilişkilidir. İyileşmeyen
  öğrencilerin tam da iyileşmedikleri için izlemi bırakması olası bir örnektir.
\end{description}

\begin{table}[htbp]
\centering
\caption{Eksik veri mekanizmalarının pratik anlamı.}
\begin{tabular}{@{}p{0.14\textwidth}p{0.30\textwidth}p{0.38\textwidth}@{}}
\toprule
Mekanizma & Eksiklik neye bağlı olabilir? & Analiz açısından temel sonuç \\
\midrule
MCAR & Ne gözlenen ne gözlenmeyen değerlere & Tam gözlem analizi yansız olabilir; hassasiyet azalır \\
MAR & Analize alınan gözlenen bilgilere & Uygun model veya çoklu atama gözlenen bilgiyi kullanabilir \\
MNAR & Gözlenmeyen değere veya kaydedilmemiş nedene & Ek varsayım ve duyarlılık analizi gerekir \\
\bottomrule
\end{tabular}
\end{table}

MCAR veriden bütünüyle doğrulanamaz. MAR ve MNAR ayrımı ise yalnız gözlenen
veriyle kesinleştirilemez; araştırma süreci, ayrılma nedenleri ve alan bilgisi
gerekir. ``Little testi anlamlı değil, dolayısıyla MCAR kanıtlandı'' biçiminde
kesin sonuç verilmemelidir.

\section{Eksik verinin sonuçları}

Eksiklik üç temel yoldan zarar verebilir:

\begin{itemize}
  \item geçerli gözlem sayısı ve test gücü azalır,
  \item standart hatalar büyür veya yanlış hesaplanır,
  \item eksik kalan kişiler sistematik farklıysa tahminler yanlılaşır.
\end{itemize}

Her analiz farklı bir gözlem altkümesi kullanırsa tablolar arasında $n$
değişebilir. Özellikle çift bazında korelasyon matrisi, her hücrede farklı
öğrencilerden hesaplanan ve birlikte geçerli bir kovaryans yapısı oluşturmayan
sonuçlar üretebilir.

\section{Tam gözlem ve mevcut gözlem yöntemleri}
\index{liste bazında silme}
\index{çift bazında silme}

\textbf{Tam gözlem analizi} veya liste bazında silme, modelde kullanılan bütün
değişkenleri gözlenen kişileri analize alır. MCAR altında yansız olabilir;
ancak bilgi kaybeder. Bazı regresyon bağlamlarında eksiklik yalnız açıklayıcı
değişkenlere ve modeldeki diğer bilgilere bağlıysa daha zayıf koşullarda da
uygun olabilir; bu sonuç her yöntem için otomatik değildir.

\textbf{Mevcut gözlem} veya çift bazında silme, her özet için o değişken
çiftinde gözlenen kayıtları kullanır. Daha çok hücreden yararlansa da analizler
farklı örneklemlere dayanır ve sonuçların birlikte yorumunu güçleştirir.

\begin{sezgibox}
Silme yöntemi ``veriyi temizlemez''; hedef evreni fiilen bütün değişkenleri
tamamlayan kişilerle sınırlar. Tamamlayanlar ile tamamlamayanlar farklıysa bu
yeni hedef, araştırmanın özgün hedefi olmayabilir.
\end{sezgibox}

\section{Tekli atama neden çoğu zaman yetersizdir?}
\index{ortalama ile atama}
\index{son gözlemi taşıma}

Ortalama, medyan, grup ortalaması veya son gözlemi ileri taşıma gibi yöntemler
her eksik hücreye tek bir değer yerleştirir. Dosya tamamlanmış görünür; fakat
atama belirsizliği hesaba katılmaz. Ortalama ile atama özellikle:

\begin{itemize}
  \item değişkenin varyansını küçültür,
  \item değişkenler arası ilişkileri zayıflatabilir,
  \item standart hata ve $p$-değerlerini gereğinden iyimser yapabilir,
  \item gözlenen dağılımı yapay bir yığılmayla bozar.
\end{itemize}

Tekli atama bazen veri giriş kontrolü, ölçek puanlama kuralı veya açıkça
tanımlanmış bir duyarlılık senaryosu için kullanılabilir; genel amaçlı çıkarım
yöntemi olarak savunulmamalıdır.

\section{Model temelli yaklaşım ve çoklu atama}
\index{çoklu atama}
\index{multiple imputation@\textit{multiple imputation}}

Çoklu atama, eksik hücreler için tek bir kesin sayı üretmez. Gözlenen veriye ve
atama modeline dayanarak $m$ farklı tamamlanmış veri seti oluşturur. Her veri
seti aynı analiz modeliyle çözümlenir; tahminler ve varyanslar Rubin birleştirme
kurallarıyla bir araya getirilir \cite{rubin1976,little2019,vanbuuren2018}.

Atama $j$ için tahmin $\hat\theta_j$ ve varyans $U_j$ olsun. Birleşik tahmin
\[
\bar\theta=\frac{1}{m}\sum_{j=1}^m\hat\theta_j
\]
olur. Atama içi ortalama varyans
\[
\bar U=\frac{1}{m}\sum_{j=1}^m U_j,
\]
atamalar arası varyans ise
\[
B=\frac{1}{m-1}\sum_{j=1}^m(\hat\theta_j-\bar\theta)^2
\]
biçimindedir. Toplam varyans
\[
T=\bar U+\left(1+\frac{1}{m}\right)B
\]
ile hem örnekleme hem atama belirsizliğini içerir.

\begin{assumptionbox}
Atama modeli; analizdeki yanıtı, açıklayıcıları, eksikliği öngören yardımcı
değişkenleri, grup ve etkileşim yapısını kapsamalıdır. Analiz modelinde bulunan
önemli bir ilişkiyi atama modelinden çıkarmak yanlı sonuç üretebilir.
\end{assumptionbox}

Çoklu atama kötü ölçümü, seçici örneklemeyi veya kaydedilmemiş bütün MNAR
nedenlerini otomatik düzeltmez. Atama sayısı, yakınsama ve dağılım tanıları,
değişken türüne uygun model ve birleştirme yöntemi ayrıca denetlenmelidir
\cite{sterne2009}.

\Needspace{30\baselineskip}
\section{Yöntem seçimi için karar akışı}

Şekil~\ref{fig:missing-data-flow}, eksik hücreyi hemen doldurmak yerine önce
veri kalitesi, hedef analiz ve mekanizma varsayımını sorgulayan akışı özetler.

\begin{figure}[H]
\centering
\begin{tikzpicture}[
  font=\small,
  decision/.style={diamond,draw=PrismBlue!80!black,fill=PrismBlue!10,
    align=center,aspect=2.25,inner sep=1pt,text width=3.1cm},
  result/.style={draw=PrismBlue!70!black,rounded corners=2pt,fill=PrismLight,
    align=center,minimum height=9mm,text width=3.2cm},
  arrow/.style={-{Stealth[length=2mm]},thick,draw=PrismBlue!80!black},
  lab/.style={font=\footnotesize,fill=white,inner sep=1pt},
  node distance=7mm and 11mm
]
\node[decision] (valid) {Değer gerçekten eksik mi?};
\node[result,right=of valid] (clean) {Kod, aralık, yinelenen kayıt veya yapısal eksikliği düzelt ve belgele};
\node[decision,below=of valid] (amount) {Eksiklik önemsiz ve MCAR varsayımı makul mü?};
\node[result,right=of amount] (complete) {Tam gözlem analizi; bilgi kaybını raporla};
\node[decision,below=of amount] (observed) {Eksikliği açıklayan gözlenen bilgiler var mı?};
\node[result,right=of observed] (mi) {Model temelli yöntem veya çoklu atama};
\node[result,below=of observed] (sensitivity) {MNAR senaryoları ve duyarlılık analizi};
\draw[arrow] (valid) -- node[lab,above]{hayır} (clean);
\draw[arrow] (valid) -- node[lab,right]{evet} (amount);
\draw[arrow] (amount) -- node[lab,above]{evet} (complete);
\draw[arrow] (amount) -- node[lab,right]{hayır} (observed);
\draw[arrow] (observed) -- node[lab,above]{evet} (mi);
\draw[arrow] (observed) -- node[lab,right]{yetersiz / hayır} (sensitivity);
\end{tikzpicture}
\caption{Eksik veri yöntemini seçmeden önce izlenecek karar akışı.}
\label{fig:missing-data-flow}
\end{figure}

Akış mekanik bir reçete değildir. Örnekleme tasarımı, boylamsal yapı,
kümeleşme, değişken türü ve bilimsel hedef seçimi değiştirebilir.

\section{Çözümlü örnek: Seçici son-test eksikliği}

İki program grubunda başlangıç ve son-test puanları Tablo~\ref{tab:missing-small-data}'de
verilmiştir. B grubundaki başlangıç puanı en düşük öğrencinin son-testi
gözlenmemiştir. Öğretim amacıyla parantez içindeki 45 değeri saklanan gerçek
değer kabul edilecek; gerçek araştırmada bu değer bilinemez.

\begin{table}[htbp]
\centering
\caption{Seçici eksiklik içeren küçük öğretim verisi.}
\label{tab:missing-small-data}
\begin{tabular}{@{}crrr@{\qquad}crrr@{}}
\toprule
No & Grup & Ön & Son & No & Grup & Ön & Son \\
\midrule
1&A&42&50 & 7&B&40&eksik (45)\\
2&A&48&55 & 8&B&44&49\\
3&A&51&59 & 9&B&49&55\\
4&A&55&62 &10&B&53&58\\
5&A&60&68 &11&B&58&64\\
6&A&63&70 &12&B&62&67\\
\bottomrule
\end{tabular}
\end{table}

A grubunun son-test ortalaması $60{,}67$'dir. B grubunda yalnız gözlenen beş
puan kullanılırsa ortalama
\[
\bar x_{B,\mathrm{gözlenen}}=58{,}60
\]
ve A--B farkı $2{,}07$ bulunur. Saklanan 45 bilinseydi B ortalaması $56{,}33$,
fark ise $4{,}33$ olacaktı. Düşük başlangıç puanlı öğrencinin kaybı B grubunu
olduğundan başarılı göstermiş ve farkı küçültmüştür.

B grubunun gözlenen ortalaması olan $58{,}60$ eksik hücreye yazılırsa B
ortalaması değişmez; fakat yapay olarak merkezde bir gözlem eklenir ve
değişkenlik küçülür. Başlangıç puanını kullanan basit regresyon tahmini ise
eksik son-test için yaklaşık $46{,}70$ üretir. Bu değer 45'e yakındır; yine de
kesin gözlem değildir ve tekli regresyon ataması tahmin belirsizliğini hesaba
katmaz.

\begin{ornek}[Tam gözlem sonucunu yorumlama]
Tam gözlem analizi ``eksik öğrencinin sonucu 58,60'tır'' demez. O öğrenciyi
hedef analizden çıkarır ve sonucu son-testi gözlenen öğrenciler için üretir.
Bu alt grubun başlangıçta farklı olması, özgün hedef evrene genellemeyi
zayıflatır.
\end{ornek}

\section{Duyarlılık analizi}
\index{duyarlılık analizi!eksik veri}

MAR varsayımı gözlenmeyen veriden doğrulanamadığı için ana analiz makul
alternatiflerle karşılaştırılmalıdır. Basit duyarlılık incelemeleri şunları
içerebilir:

\begin{itemize}
  \item tam gözlem ve çoklu atama sonuçlarını karşılaştırmak,
  \item atanan son-test değerlerini belirli miktarda daha düşük veya yüksek almak,
  \item kayıp öğrenciler için en iyi ve en kötü makul senaryoları kurmak,
  \item farklı yardımcı değişken kümeleriyle atama yapmak,
  \item sonuç yönü, büyüklüğü ve kararının hangi varsayımlarda değiştiğini göstermek.
\end{itemize}

Amaç en uygun sonucu seçmek değil, sonucun hangi varsayıma ne kadar bağımlı
olduğunu görünür kılmaktır.

\section{Python ile veri kalite denetimi}

\begin{lstlisting}[style=prismPython]
import numpy as np
import pandas as pd

df = pd.DataFrame({
    "grup": ["A"]*6 + ["B"]*6,
    "on":   [42,48,51,55,60,63,40,44,49,53,58,62],
    "son":  [50,55,59,62,68,70,np.nan,49,55,58,64,67]
})

# Değişken ve kişi düzeyinde eksiklik
print(df.isna().sum())
print(df.isna().mean())
df["eksik_son"] = df["son"].isna().astype(int)
print(pd.crosstab(df["grup"], df["eksik_son"], margins=True))

# Gözlenen ve eksik son-test gruplarının başlangıç puanları
print(df.groupby("eksik_son")["on"].agg(["count", "mean", "std"]))

# Tam gözlem özeti
tam = df.dropna(subset=["son"])
print(tam.groupby("grup")["son"]
         .agg(["count", "mean", "std"]))
\end{lstlisting}

Kod, eksiklik göstergesinin başlangıç puanıyla ilişkisini betimler; tek başına
MCAR, MAR veya MNAR kanıtı değildir. Çoklu atama gerekiyorsa Python'da
\texttt{statsmodels.imputation.mice} gibi bir yordam kullanılabilir; atama
modeli analiz modeliyle uyumlu kurulmalıdır.

\section{R ile veri kalite denetimi ve çoklu atama iskeleti}

\begin{lstlisting}[style=prismR]
veri <- data.frame(
  grup = factor(rep(c("A", "B"), each = 6)),
  on  = c(42,48,51,55,60,63,40,44,49,53,58,62),
  son = c(50,55,59,62,68,70,NA,49,55,58,64,67)
)

colSums(is.na(veri))
colMeans(is.na(veri))
veri$eksik_son <- as.integer(is.na(veri$son))
with(veri, table(grup, eksik_son, useNA = "ifany"))
aggregate(on ~ eksik_son, data = veri,
          FUN = function(x) c(n=length(x), ort=mean(x), ss=sd(x)))

complete <- subset(veri, !is.na(son))
aggregate(son ~ grup, data = complete,
          FUN = function(x) c(n=length(x), ort=mean(x), ss=sd(x)))

# mice paketi kuruluysa örnek çoklu atama akışı:
# imp <- mice::mice(veri[c("grup","on","son")], m=20,
#                   method="pmm", seed=2026, printFlag=FALSE)
# fit <- with(imp, lm(son ~ on + grup))
# summary(mice::pool(fit), conf.int=TRUE)
\end{lstlisting}

\texttt{mice} örneğindeki öngörücü ortalama eşleştirmesi, gözlenen değerlere
benzeyen atamalar üretmeyi amaçlar. Paket adı yazmak yöntem gerekçesinin yerine
geçmez; değişken türleri, etkileşimler ve tanılar ayrıca incelenir.

\section{SPSS ile eksik veri denetimi}

Önce özel eksik kodlar tanımlanmalı veya sistem eksik değere dönüştürülmelidir.
Aşağıdaki temel sözdizimi eksiklik göstergesi ve grup tablosu üretir:

\begin{lstlisting}[style=prismSPSS]
MISSING VALUES son (99, 999).
COMPUTE eksik_son = MISSING(son).
VALUE LABELS eksik_son 0 'Gözlendi' 1 'Eksik'.
EXECUTE.

FREQUENCIES VARIABLES=on son eksik_son.
CROSSTABS
 /TABLES=grup BY eksik_son
 /CELLS=COUNT ROW.

MEANS TABLES=on BY eksik_son
 /CELLS=COUNT MEAN STDDEV.
\end{lstlisting}

Menüde temel inceleme için \textbf{Analyze $>$ Descriptive Statistics}, daha
ayrıntılı örüntü incelemesi için kurulumda varsa \textbf{Analyze $>$ Missing
Value Analysis} kullanılabilir. Çoklu atama modülü bulunan sürümlerde
\textbf{Analyze $>$ Multiple Imputation $>$ Impute Missing Data Values}
yoluyla atama modeli kurulur. Üretilen çoklu dosyalarda analiz sonuçlarının
birleştirilmiş tablosu kullanılmalı; tek bir atama dosyasının sonucu nihai
sonuç gibi raporlanmamalıdır.

\begin{outputguidebox}
Çıktıyı şu sırayla okuyun: (1) özel eksik kodlar doğru tanımlanmış mı,
(2) her değişken için geçerli ve eksik $n$, (3) kişi ve örüntü düzeyindeki
eksiklik, (4) gruplara göre eksiklik oranı, (5) eksikliği öngören gözlenen
değişkenler, (6) kullanılan silme/atama modeli, (7) birleşik tahmin ve standart
hata, (8) duyarlılık sonucu.
\end{outputguidebox}

\section{Bilimsel raporlama örneği}

``Son-test puanı 120 öğrencinin 102'sinde gözlenmiş, 18 öğrencide (yüzde 15)
eksik kalmıştır. Eksiklik program grubu ve başlangıç kaygı puanına göre
farklılaştığından tamamen rassal eksiklik varsayımı uygun görülmemiştir. Ana
analizde grup, başlangıç puanı ve devam bilgilerini içeren 20 veri setli çoklu
atama kullanılmış; tahminler Rubin kurallarıyla birleştirilmiştir. Tam gözlem
analizi ile son-test puanlarının 3 puan daha düşük olduğu MNAR duyarlılık
senaryosu ek olarak raporlanmıştır. Sonucun yönü korunmuş, büyüklüğü duyarlılık
senaryosunda azalmıştır.''

\section{Bir eksik veri analizini değerlendirirken}

\begin{enumerate}
  \item Eksik, sıfır, uygulanamaz ve geçersiz kodlar ayrılmış mı?
  \item Değişken, kişi, grup ve örüntü düzeyindeki eksiklik raporlanmış mı?
  \item Mekanizma varsayımı araştırma süreci ve gözlenen bilgilerle gerekçelendirilmiş mi?
  \item Silme veya atama yöntemi hedef analizle uyumlu mu?
  \item Atama modeli sonuç, açıklayıcı ve yardımcı değişkenleri içeriyor mu?
  \item Atama belirsizliği standart hataya yansıtılmış mı?
  \item Tam gözlem ve MNAR duyarlılık sonuçları karşılaştırılmış mı?
  \item Temizleme ve atama kodu yeniden üretilebilir biçimde saklanmış mı?
\end{enumerate}

\section{Bölüm sonu çözümlü problemler}

\begin{enumerate}
  \item \textbf{Eksiklik oranı.} 240 öğrencinin 18'inde başlangıç puanı,
  36'sında son-test puanı eksiktir. Değişken bazındaki eksiklik oranlarını
  bulun.

  \textbf{Çözüm:} Başlangıç için $18/240=0{,}075$, yani yüzde 7,5; son-test
  için $36/240=0{,}15$, yani yüzde 15'tir. Aynı öğrencilerin iki değişkende de
  eksik olup olmadığı bilinmeden tam gözlem sayısı yalnız bu oranlardan
  bulunamaz.

  \item \textbf{Mekanizma yorumu.} Son-test eksikliği yalnız gözlenen
  başlangıç puanı ve program grubuyla ilişkili görünmektedir. Hangi mekanizma
  çalışma varsayımı olarak düşünülebilir?

  \textbf{Çözüm:} Başlangıç puanı ve grup atama/analiz modeline alındığında MAR
  makul bir çalışma varsayımı olabilir. Gözlenmeyen son-test değerinin ayrıca
  eksikliği etkileyip etkilemediği gözlenen veriden kanıtlanamaz; MNAR
  duyarlılık analizi gerekir.

  \item \textbf{Ortalama ile atama.} Gözlenen değerler $40,50,60$ ve bir eksik
  hücre olsun. Eksik hücreye gözlenen ortalama yazıldığında ortalama ve
  örneklem varyansı nasıl değişir?

  \textbf{Çözüm:} Gözlenen ortalama 50'dir. $40,50,60,50$ dizisinin ortalaması
  yine 50 olur. İlk üç değerin örneklem varyansı 100, atama sonrasındaki
  örneklem varyansı $200/3\approx66{,}67$ olur. Yapay merkez değeri değişkenliği
  küçültür.

  \item \textbf{Rubin toplam varyansı.} $m=5$, atama içi ortalama varyans
  $\bar U=4$ ve atamalar arası varyans $B=1{,}5$ olsun. Toplam varyansı ve
  yaklaşık standart hatayı bulun.

  \textbf{Çözüm:}
  \[
  T=4+\left(1+\frac15\right)(1{,}5)=5{,}8,
  \qquad \SE=\sqrt{5{,}8}\approx2{,}41.
  \]
  Yalnız $\sqrt{\bar U}=2$ kullanmak atama belirsizliğini dışarıda bırakır.

  \item \textbf{Duyarlılık kararı.} MAR çoklu atama analizinde grup farkı 5,2
  puan; eksik sonuçların her grupta 4 puan daha düşük varsayıldığı senaryoda
  4,9 puan; yalnız müdahale grubunda 8 puan daha düşük varsayıldığı senaryoda
  1,1 puandır. Sonucu değerlendirin.

  \textbf{Çözüm:} Benzer düşüş varsayımında sonuç kararlıdır. Müdahale
  grubundaki güçlü seçici kayıp varsayımında fark büyük ölçüde azalmaktadır.
  Ana bulgu eksikliğin gruplar arasında farklı davranmadığı varsayımına
  duyarlıdır; bu sınır açıkça raporlanmalıdır.
\end{enumerate}

\bolumkaynaklari{\cite{rubin1976,little2019,vanbuuren2018,sterne2009}}

\section*{Hatırla--Uygula--Yansıt}

\begin{enumerate}
\renewcommand{\labelenumi}{\textbf{\Alph{enumi}.}}
  \item \textbf{Ayır:} Boş hücre, gerçek sıfır, 99 kodu ve yapısal eksiklik için farklı işlem yazın.
  \item \textbf{Özetle:} 200 kişilik veride bir değişkende 30 eksik varsa oranı bulun; neden genel yüzdeyle yetinilmemesi gerektiğini açıklayın.
  \item \textbf{Mekanizma:} MCAR, MAR ve MNAR için PDR bağlamında birer örnek verin.
  \item \textbf{Eleştir:} ``Eksiklik yüzde 5'ten azsa silmek her zaman güvenlidir'' ifadesini düzeltin.
  \item \textbf{Karşılaştır:} Liste bazında ve çift bazında silmenin temel farkını yazın.
  \item \textbf{Varyans:} Ortalama ile atamanın neden değişkenliği küçülttüğünü açıklayın.
  \item \textbf{Birleştir:} Çoklu atamada atama içi ve atamalar arası varyansın rollerini ayırın.
  \item \textbf{Planla:} Son-test eksikliği için atama modeline alınabilecek dört gözlenen değişken önerin.
  \item \textbf{Duyarlılık:} MNAR olasılığı için uygulanabilecek bir senaryo analizi tasarlayın.
  \item \textbf{Raporla:} Eksiklik miktarı, yöntem, varsayım ve duyarlılık sonucunu içeren kısa yöntem paragrafı yazın.
\end{enumerate}